\documentclass[conference]{IEEEtran}
\IEEEoverridecommandlockouts

\usepackage{cite}
\usepackage{amsmath,amssymb,amsfonts}
\usepackage{algorithmic}
\usepackage{algorithm}
\usepackage{graphicx}
\usepackage{textcomp}
\usepackage{xcolor}
\usepackage{booktabs}
\usepackage{multirow}
\usepackage{url}
\usepackage{hyperref}
\usepackage{listings}
\usepackage{array}
\usepackage{tabularx}

\hypersetup{
    colorlinks=true,
    linkcolor=blue,
    citecolor=blue,
    urlcolor=blue
}

\lstset{
    basicstyle=\ttfamily\footnotesize,
    breaklines=true,
    frame=single,
    numbers=left,
    numberstyle=\tiny,
    xleftmargin=2em,
}

\def\BibTeX{{\rm B\kern-.05em{\sc i\kern-.025em b}\kern-.08em
    T\kern-.1667em\lower.7ex\hbox{E}\kern-.125emX}}

\begin{document}

\title{LBRO: A Law-Aware Breach Response Orchestrator for Automated Cyber Incident Detection, Evidence Preservation, and Multi-Jurisdiction Compliance Enforcement}

\author{
\IEEEauthorblockN{[Author 1]}
\IEEEauthorblockA{\textit{Department of [Department]} \\
\textit{[University/Institution]}\\
[City, Country] \\
[email@institution.edu]}
\and
\IEEEauthorblockN{[Author 2]}
\IEEEauthorblockA{\textit{Department of [Department]} \\
\textit{[University/Institution]}\\
[City, Country] \\
[email@institution.edu]}
\and
\IEEEauthorblockN{[Author 3]}
\IEEEauthorblockA{\textit{Department of [Department]} \\
\textit{[University/Institution]}\\
[City, Country] \\
[email@institution.edu]}
}

\maketitle

% ─────────────────────────────────────────────────────────────────────────────
\begin{abstract}
The escalating frequency and sophistication of cyberattacks have created a critical gap in incident response capabilities, particularly for small and medium-sized organizations lacking dedicated Security Operations Centers (SOC). Existing solutions are either prohibitively expensive, fragmented across multiple tools, or require significant security expertise to operate. This paper presents LBRO (Law-aware Breach Response Orchestrator), an integrated, open-architecture platform that unifies automated threat detection, machine learning (ML)-based attack classification, tamper-evident evidence preservation, and multi-jurisdiction regulatory compliance in a single deployable system. LBRO employs a Gaussian Naive Bayes classifier trained on a stratified 15,000-sample subset of the CICIDS2017 dataset spanning 77 network flow features and 15 attack categories, achieving a macro F1 score of 0.9692, balanced accuracy of 0.9678, Matthews Correlation Coefficient (MCC) of 0.9957, and a 3-fold cross-validation F1 of 0.9595. The backend, implemented in Python with FastAPI and asynchronous PostgreSQL via SQLAlchemy 2.0, exposes a versioned REST API secured by JSON Web Tokens (JWT) with per-token revocation, bcrypt password hashing, and a four-tier Role-Based Access Control (RBAC) system enforcing 30 granular permissions. The Evidence Vault stores forensic artifacts with SHA-256 integrity hashing, immutability flags, and tamper-evident chain-of-custody logging for every access and download event. A rule-based Compliance Engine automatically generates binding obligations under GDPR (72-hour notification), HIPAA (60-day disclosure), and India's DPDPA (72-hour notification) based on incident jurisdiction and data classification. The frontend, built with React 18 and TypeScript, provides a real-time seven-tab investigation workspace, Server-Sent Events (SSE) live stream, and automated PDF report generation. The full stack is containerized with Docker Compose and includes Terraform Infrastructure-as-Code for production deployment on AWS ECS. Experimental evaluation demonstrates sub-millisecond ML inference latency, complete API response times under 250 ms, and a 19 KB model footprint suitable for resource-constrained deployments. LBRO provides a reproducible, law-aware foundation for incident response research and operational deployment.
\end{abstract}

\begin{IEEEkeywords}
Cybersecurity, Incident Response, Network Intrusion Detection, Machine Learning, CICIDS2017, Gaussian Naive Bayes, RBAC, Digital Forensics, Evidence Preservation, GDPR, HIPAA, DPDPA, FastAPI, Docker, MITRE ATT\&CK, Compliance Automation, Security Orchestration
\end{IEEEkeywords}

% ─────────────────────────────────────────────────────────────────────────────
\section{Introduction}

The global cost of cybercrime is projected to reach \$10.5 trillion annually by 2025 \cite{morgan2020cybercrime}, yet the tools available to organizations for detecting and responding to breaches remain fragmented, expensive, and difficult to operate without specialized staff. A typical enterprise incident response workflow spans multiple disconnected products: a SIEM for log aggregation, a SOAR platform for automation, a forensic tool for evidence collection, and separate legal counsel to navigate regulatory obligations. For small and medium enterprises (SMEs)—which constitute over 90\% of businesses globally \cite{oecd2021}—this fragmentation is operationally impractical.

The problem is compounded by evolving legal obligations. The European Union's General Data Protection Regulation (GDPR) mandates notification of data protection authorities within 72 hours of breach discovery \cite{voigt2017eu}. The United States Health Insurance Portability and Accountability Act (HIPAA) requires disclosure to the Department of Health and Human Services within 60 days \cite{annas2003hipaa}. India's Digital Personal Data Protection Act (DPDPA) 2023 similarly imposes a 72-hour notification deadline \cite{dpdpa2023}. Non-compliance carries fines up to 4\% of global annual turnover under GDPR and up to INR 250 crore under DPDPA. Organizations lack tooling that integrates legal awareness into the technical incident response workflow.

Machine learning has emerged as a promising approach for automated threat detection. The CICIDS2017 dataset \cite{sharafaldin2018cicids}, generated at the Canadian Institute for Cybersecurity, provides labeled network flow records across 15 attack categories and has become a de facto benchmark for intrusion detection research \cite{panigrahi2018survey}. However, most published work evaluates classifiers in isolation without integrating them into operational incident response pipelines.

This paper makes the following contributions:

\begin{enumerate}
    \item \textbf{LBRO platform}: A complete, production-grade incident response system integrating ML classification, evidence preservation, and compliance enforcement in a single deployable stack.
    \item \textbf{ML evaluation}: A systematic comparison of nine classifiers on CICIDS2017 features using a composite scoring function that balances macro F1, balanced accuracy, MCC, and cross-validation stability; identifying Gaussian Naive Bayes as the optimal deployment choice with a composite score of 0.9731.
    \item \textbf{Law-aware compliance engine}: An automated obligation generator that maps incident attributes to GDPR, HIPAA, and DPDPA requirements with deadline tracking and persistent state management.
    \item \textbf{Tamper-evident forensics}: A chain-of-custody evidence model with SHA-256 integrity verification, immutability enforcement, and per-access audit records.
    \item \textbf{Project-scoped ingestion API}: A multi-tenant event ingestion architecture using per-project cryptographic API keys that enable zero-configuration external application monitoring.
\end{enumerate}

The remainder of this paper is organized as follows: Section II reviews related work; Section III defines the problem; Section IV describes the system architecture; Section V details the implementation; Section VI presents the ML component; Section VII reports experimental results; Section VIII analyzes the security model; Section IX discusses the compliance framework; Section X discusses findings and limitations; and Section XI concludes.

% ─────────────────────────────────────────────────────────────────────────────
\section{Related Work}

\subsection{Commercial SIEM and SOAR Platforms}

Splunk Enterprise Security \cite{splunk2023} is the dominant commercial SIEM, providing log aggregation, correlation rules, and dashboards. However, it requires significant infrastructure investment (typically \$50,000+ annually for enterprise licenses), specialized Splunk Query Language (SPL) expertise, and does not natively integrate legal compliance or evidence chain-of-custody. IBM QRadar \cite{ibm2023qradar} offers similar capabilities with automated threat intelligence feeds but shares comparable cost and operational complexity barriers.

CrowdStrike Falcon \cite{crowdstrike2022} provides cloud-native endpoint detection and response (EDR), while Microsoft Defender for Endpoint \cite{microsoft2023defender} integrates with the Azure ecosystem. Both focus primarily on endpoint telemetry rather than application-layer security events and do not provide regulatory compliance automation.

Elastic Security \cite{elastic2023} offers an open-source path through the ELK stack (Elasticsearch, Logstash, Kibana) with a security solution layer. While cost-effective, it requires substantial engineering effort to configure, lacks built-in compliance frameworks, and does not provide evidence preservation with chain-of-custody.

\subsection{Open-Source Security Platforms}

Wazuh \cite{wazuh2023} is an open-source SIEM and XDR platform built on the ELK stack. It provides log analysis, vulnerability detection, and compliance reporting for PCI-DSS and HIPAA. Wazuh requires agent installation on monitored endpoints and does not provide ML-based attack classification or automated incident investigation workspaces.

TheHive \cite{thehive2023} is an open-source Security Incident Response Platform (SIRP) that provides case management and observable analysis. It requires integration with Cortex analyzers for ML enrichment and does not include built-in network traffic classification or compliance deadline management. Shuffle SOAR \cite{shuffle2022} provides workflow automation for security operations but lacks native ML classification, evidence storage, or compliance engines.

\subsection{Network Intrusion Detection Research}

Buczak and Guven \cite{buczak2016survey} conducted a comprehensive survey of ML methods for network intrusion detection, evaluating decision trees, na\"{i}ve Bayes, neural networks, and ensemble methods. They identify the trade-off between detection accuracy and computational cost as the primary deployment challenge. Liao et al. \cite{liao2013ids} reviewed IDS architectures and noted that most research systems fail to address operational deployment concerns such as model update cycles, latency constraints, and integration with existing workflows.

Sharafaldin et al. \cite{sharafaldin2018cicids} introduced the CICIDS2017 dataset, providing 78 features derived from network traffic using CICFlowMeter. Subsequent research by Panigrahi and Borah \cite{panigrahi2018survey} systematically evaluated 37 ML methods on CICIDS2017, finding that ensemble methods (Random Forest, XGBoost) achieve the highest raw accuracy but at significantly higher computational cost than probabilistic classifiers.

\subsection{Digital Forensics and Evidence Integrity}

Casey \cite{casey2011digital} established the foundational principles of digital forensics evidence handling, emphasizing the chain of custody as the mechanism that establishes the admissibility and integrity of digital evidence. Garfinkel \cite{garfinkel2010dfr} identified evidence integrity verification (hash comparison across storage events) and structured audit trails as the two minimum requirements for legally defensible digital forensics. LBRO implements both through its SHA-256 evidence hashing and the \texttt{chain\_of\_custody} table that records every access, download, and verification event.

\subsection{Regulatory Compliance and Automation}

Goddard \cite{goddard2017gdpr} analyzed the GDPR's technical implications, identifying automated breach detection and rapid notification as the primary engineering challenges for compliance. Mercuri \cite{mercuri2004hipaa} examined HIPAA's technical safeguard requirements and found that audit logging and access control are the most commonly violated provisions. LBRO addresses both through its RBAC system and immutable audit trail.

\subsection{Gap Analysis}

Table~\ref{tab:comparison} summarizes the feature comparison between LBRO and existing platforms. No existing open-source platform integrates ML-based network traffic classification, tamper-evident evidence preservation, multi-jurisdiction compliance automation, and a project-scoped ingestion API in a single deployable unit.

\begin{table*}[t]
\caption{Feature Comparison: LBRO vs. Existing Security Platforms}
\label{tab:comparison}
\centering
\begin{tabular}{@{}lccccccccc@{}}
\toprule
\textbf{Feature} & \textbf{LBRO} & \textbf{Splunk ES} & \textbf{Wazuh} & \textbf{CrowdStrike} & \textbf{Elastic Sec.} & \textbf{QRadar} & \textbf{TheHive} & \textbf{Shuffle} \\
\midrule
ML Attack Classification   & \checkmark & Partial & -- & \checkmark & Partial & \checkmark & Via plugin & -- \\
Evidence Chain of Custody  & \checkmark & -- & -- & -- & -- & -- & \checkmark & -- \\
GDPR Compliance Engine     & \checkmark & Partial & Partial & -- & -- & Partial & -- & -- \\
HIPAA Compliance Engine    & \checkmark & \checkmark & \checkmark & -- & -- & \checkmark & -- & -- \\
DPDPA Compliance Engine    & \checkmark & -- & -- & -- & -- & -- & -- & -- \\
Project API Key Ingestion  & \checkmark & -- & -- & -- & -- & -- & -- & -- \\
RBAC (4 roles)             & \checkmark & \checkmark & \checkmark & \checkmark & \checkmark & \checkmark & \checkmark & Partial \\
Open Source / Self-hosted  & \checkmark & -- & \checkmark & -- & \checkmark & -- & \checkmark & \checkmark \\
Docker Compose Deploy      & \checkmark & -- & \checkmark & -- & \checkmark & -- & \checkmark & \checkmark \\
PDF Report Generation      & \checkmark & \checkmark & -- & \checkmark & -- & \checkmark & -- & -- \\
Live Event Stream (SSE)    & \checkmark & -- & -- & -- & -- & -- & -- & -- \\
Automated Incident Creation& \checkmark & Partial & Partial & \checkmark & Partial & \checkmark & -- & Via workflow \\
\bottomrule
\end{tabular}
\end{table*}

% ─────────────────────────────────────────────────────────────────────────────
\section{Problem Statement}

Current incident response tooling exhibits four primary limitations that motivate LBRO's design:

\textbf{P1 — Manual and fragmented response.} Incident response is distributed across separate detection, ticketing, forensic, and reporting tools. The median time to identify and contain a breach is 277 days \cite{ibm2022cost}, partly attributable to this fragmentation. LBRO addresses this by providing a unified investigation workspace with ML classification, evidence management, timeline reconstruction, and IOC extraction in a single interface.

\textbf{P2 — Compliance opacity.} Organizations do not know which regulations apply to a given breach until legal review is complete, introducing delays that cause non-compliance by default. LBRO's Compliance Engine evaluates each incident against GDPR, HIPAA, and DPDPA criteria at creation time, immediately surfacing applicable obligations with countdown deadlines.

\textbf{P3 — Evidence integrity risk.} Digital evidence collected without a strict chain of custody may be inadmissible or legally contested. File copies can be silently corrupted; access is rarely logged. LBRO computes SHA-256 hashes at upload time, enforces immutability at the database layer (\texttt{is\_immutable=True}), and records every access, download, and hash-verification event in the \texttt{chain\_of\_custody} table with a timestamp, actor identity, and IP address.

\textbf{P4 — Integration barrier.} Monitored applications must undergo significant instrumentation to integrate with existing SIEMs (agent installation, log forwarding configuration, schema mapping). LBRO provides a single HTTP POST endpoint authenticated by a project-scoped API key, requiring as little as three lines of application code to begin ingesting security events.

% ─────────────────────────────────────────────────────────────────────────────
\section{System Architecture}

\subsection{High-Level Design}

LBRO follows a layered microservice architecture comprising five logical tiers: (1) presentation, (2) API gateway, (3) business logic, (4) data persistence, and (5) background processing. Fig.~\ref{fig:architecture} illustrates the overall architecture.

\begin{figure}[h]
\centering
\fbox{\parbox{0.9\columnwidth}{
\centering
\vspace{4pt}
\textbf{LBRO System Architecture}\\[4pt]
\small
\texttt{[React/TS Frontend]} $\rightarrow$ \texttt{[Nginx Reverse Proxy]}\\
$\downarrow$\\
\texttt{[FastAPI REST API + SSE]} \\
$\downarrow$ (JWT + RBAC + Project Key) \\
\texttt{[Routers: Auth | Incidents | Evidence]}\\
\texttt{[Compliance | ML | Events | Reports]}\\
$\downarrow$ \\
\texttt{[SQLAlchemy Async]} $\rightarrow$ \texttt{[PostgreSQL 16]}\\
$\downarrow$ \\
\texttt{[ML Classifier]} \texttt{[Evidence Vault]} \texttt{[Compliance Engine]}\\
$\downarrow$ \\
\texttt{[AWS S3 / SQS / SecretsManager (LocalStack dev)]}
\vspace{4pt}
}}
\caption{LBRO layered system architecture.}
\label{fig:architecture}
\end{figure}

\subsection{Authentication and Authorization}

LBRO implements a dual authentication model. User sessions use JWT access tokens (HS256, 30-minute expiry) paired with refresh tokens (7-day expiry). Each JWT carries a unique \texttt{jti} claim enabling per-token revocation on logout without invalidating all sessions. Password hashing uses bcrypt via \texttt{passlib}, with account lockout enforced after 5 consecutive failed attempts for 15 minutes.

External application ingestion uses project-scoped API keys of the form \texttt{proj\_<cryptographic\_token>}, generated via Python's \texttt{secrets.token\_urlsafe(32)}. The project identity is resolved server-side from the API key; clients cannot specify or forge a project identifier.

\subsection{Role-Based Access Control}

LBRO implements a two-level RBAC hierarchy:

\begin{itemize}
    \item \textbf{Platform level}: \texttt{super\_admin} — manages the entire deployment, bypasses project boundaries for administrative operations. Platform bypass events are always audit-logged.
    \item \textbf{Project level}: \texttt{admin} $>$ \texttt{analyst} $>$ \texttt{viewer} — scoped to a single project.
\end{itemize}

Authorization decisions are made exclusively through a centralized \texttt{ROLE\_PERMISSIONS} dictionary mapping roles to 30 enumerated \texttt{Permission} values (e.g., \texttt{incident:create}, \texttt{evidence:upload}, \texttt{compliance:manage}). Role strings are never compared directly in application code; this ensures consistent enforcement across all 16 API routers.

\subsection{Project Isolation}

Each \texttt{Project} record holds a unique cryptographic API key and acts as a data boundary. All incidents, security events, evidence, compliance records, and compliance obligations are scoped to a project via foreign key relationships. Multi-tenant isolation is enforced at the ORM query layer: every data retrieval operation filters by the authenticated user's project context. This prevents cross-project data access regardless of authentication state.

\subsection{Event Ingestion Pipeline}

The ingestion pipeline, exposed at \texttt{POST /api/v1/events} and \texttt{POST /api/v1/events/batch}, processes security events in six stages:

\begin{enumerate}
    \item \textbf{API key authentication}: Resolve project from \texttt{Authorization: Bearer proj\_*} header.
    \item \textbf{Schema validation}: Pydantic model validation; invalid events return HTTP 422.
    \item \textbf{Persistence}: \texttt{SecurityEvent} written to PostgreSQL with \texttt{processing\_status=pending}.
    \item \textbf{ML classification}: CICIDS2017 GNB classifier assigns \texttt{ml\_attack\_category} and \texttt{ml\_confidence}.
    \item \textbf{Incident auto-creation}: Events with \texttt{severity \in \{high, critical\}} trigger automatic \texttt{Incident} record creation.
    \item \textbf{SSE broadcast}: Processed event is published to the Server-Sent Events bus for real-time frontend delivery.
\end{enumerate}

The batch endpoint accepts up to 1,000 events per request with partial success semantics; invalid events are skipped and the response includes per-category accepted/rejected counts.

\subsection{Database Schema}

The database comprises 14 tables managed by 11 Alembic migrations. Core entities are:

\begin{itemize}
    \item \texttt{users} — authentication, lockout, RBAC role
    \item \texttt{incidents} — security incidents with ML attributes, network features (JSON), jurisdiction flags
    \item \texttt{incident\_actions} — timeline of response actions per incident
    \item \texttt{evidence} — forensic artifacts with SHA-256 hash, deferred \texttt{file\_data} (LargeBinary), immutability flag
    \item \texttt{chain\_of\_custody} — access audit trail per evidence record
    \item \texttt{compliance\_records} — per-incident regulatory obligations with deadlines
    \item \texttt{compliance\_obligations} — project-level control tracking (GDPR/HIPAA/DPDPA articles)
    \item \texttt{compliance\_assessments} — point-in-time compliance snapshots
    \item \texttt{projects} — multi-tenant isolation unit with API key
    \item \texttt{security\_events} — raw events from external ingestion
    \item \texttt{investigation\_notes} — analyst notes with cascade delete
    \item \texttt{revoked\_tokens} — JTI blacklist for logout revocation
    \item \texttt{audit\_logs} — immutable system event log
    \item \texttt{notifications} — compliance and alert notifications
\end{itemize}

All primary keys are UUID v4. Foreign keys use \texttt{ON DELETE CASCADE} or \texttt{SET NULL} semantics appropriate to the entity relationship. Performance-critical columns (severity, status, project\_id, source\_ip, incident\_id) carry B-tree indexes added in migration 009.

% ─────────────────────────────────────────────────────────────────────────────
\section{Implementation}

\subsection{Backend}

The backend is implemented in Python 3.12 using FastAPI 0.115.5 with Uvicorn as the ASGI server. Database access uses SQLAlchemy 2.0 ORM in fully asynchronous mode via the \texttt{asyncpg} driver and \texttt{async\_sessionmaker} factory. The complete dependency set is declared in \texttt{requirements.txt}; key libraries are listed in Table~\ref{tab:deps}.

\begin{table}[h]
\caption{Key Backend Dependencies}
\label{tab:deps}
\centering
\begin{tabular}{@{}lll@{}}
\toprule
\textbf{Library} & \textbf{Version} & \textbf{Purpose} \\
\midrule
FastAPI       & 0.115.5  & REST API framework \\
SQLAlchemy    & 2.0.36   & Async ORM \\
asyncpg       & 0.30.0   & PostgreSQL async driver \\
Alembic       & 1.14.0   & Database migrations \\
python-jose   & 3.3.0    & JWT creation/verification \\
passlib       & 1.7.4    & bcrypt password hashing \\
scikit-learn  & 1.5.2    & ML training and inference \\
numpy         & 1.26.4   & Feature vector manipulation \\
reportlab     & 4.2.5    & PDF report generation \\
structlog     & 26.1.0   & Structured JSON logging \\
boto3         & 1.35.74  & AWS S3/SQS integration \\
pydantic      & 2.10.3   & Schema validation/settings \\
\bottomrule
\end{tabular}
\end{table}

The application exposes 16 API routers mounted under the \texttt{/api/v1} prefix, covering authentication, incidents, evidence, compliance, users, ML insights, dashboard, audit, reports, projects, platform management, event ingestion, demo data generation, and notifications.

Middleware is applied in three layers: (1) \texttt{SecurityHeadersMiddleware} appends \texttt{X-Frame-Options: DENY}, \texttt{X-Content-Type-Options: nosniff}, \texttt{X-XSS-Protection: 1; mode=block}, \texttt{Referrer-Policy: strict-origin-when-cross-origin}, and environment-appropriate Content Security Policy headers; (2) \texttt{RateLimitMiddleware} implements an in-memory sliding-window limiter with 60 requests/minute for general endpoints and strict 10 requests/minute for authentication endpoints; (3) \texttt{TrustedHostMiddleware} restricts the \texttt{Host} header to known values.

\subsection{Frontend}

The frontend is a single-page application built with React 18.2 and TypeScript 5.3, bundled by Vite 5. State management uses Zustand for global auth/project state and TanStack Query v5 for server state with automatic cache invalidation. UI components use Radix UI primitives with Tailwind CSS and custom styling.

The incident investigation workspace provides seven tabs: (1) Overview — key identifiers, network context, security classification; (2) Timeline — chronological event reconstruction from incident lifecycle; (3) Raw Event — raw JSON record with expand/collapse and download; (4) Evidence — vault with authenticated download, preview, hash verification, and chain of custody; (5) Notes — analyst investigation notes with inline editing; (6) IOC — extracted indicators (IPs, ports, SHA-256 hashes, protocols) with CSV export; (7) Related — incidents sharing source IP, attack category, project, or analyst.

\subsection{Containerization and Deployment}

The development stack is orchestrated by Docker Compose with six services: \texttt{postgres} (PostgreSQL 16 Alpine), \texttt{localstack} (AWS service emulation for S3, SQS, and Secrets Manager), \texttt{migrate} (one-shot Alembic migration runner), \texttt{api} (FastAPI application), \texttt{worker} (background incident/notification processor), and \texttt{frontend} (React/Nginx). All services include health checks and restart policies. A \texttt{seed} service creates default administrative users on first startup.

For production, Terraform Infrastructure-as-Code modules provision AWS ECS (Fargate) clusters, RDS PostgreSQL instances, S3 buckets for evidence and reports, SQS queues for asynchronous processing, WAF rules, VPC with private subnets, and AWS Secrets Manager for credential management. SDK clients in Python, Node.js, Go, and Java are provided for external application integration.

\subsection{PDF Report Generation}

Incident reports are generated on demand via \texttt{GET /api/v1/incidents/\{id\}/report} and returned as streaming PDF responses with \texttt{Content-Disposition: attachment}. The reportlab 4.2.5 library produces multi-page documents including: cover page with classification header, executive summary, business and technical impact analysis, network context table, ML classification with confidence score, MITRE ATT\&CK and OWASP mapping, evidence table with SHA-256 hashes, remediation recommendations, and compliance notification status table.

% ─────────────────────────────────────────────────────────────────────────────
\section{Machine Learning Component}

\subsection{Dataset}

The ML classifier is trained on a 15,000-sample stratified subset of the CICIDS2017 \cite{sharafaldin2018cicids} dataset. The dataset contains labeled network flow records derived from five days of realistic traffic generation at the Canadian Institute for Cybersecurity. Table~\ref{tab:dataset} summarizes the dataset properties. The class distribution is moderately imbalanced, with BENIGN comprising approximately 51\% of samples and rare classes such as Heartbleed and Infiltration below 0.5\%.

\begin{table}[h]
\caption{CICIDS2017 Dataset Properties (Stratified Subset)}
\label{tab:dataset}
\centering
\begin{tabular}{@{}ll@{}}
\toprule
\textbf{Property} & \textbf{Value} \\
\midrule
Total samples & 15,000 \\
Training samples & 12,000 (80\%) \\
Test samples & 3,000 (20\%) \\
Feature count & 77 \\
Attack classes & 15 \\
BENIGN proportion & $\approx$51\% \\
Rare class proportion & $<$0.5\% \\
Split strategy & Stratified random \\
\bottomrule
\end{tabular}
\end{table}

The 15 attack categories are: BENIGN, DoS Hulk, PortScan, DDoS, DoS GoldenEye, FTP-Patator, SSH-Patator, DoS slowloris, DoS Slowhttptest, Bot, Web Attack -- Brute Force, Web Attack -- XSS, Infiltration, Web Attack -- SQL Injection, and Heartbleed.

\subsection{Feature Engineering}

All 77 CICIDS2017 canonical features are used, covering four categories of network flow statistics: (1) packet-level statistics (lengths, inter-arrival times, flag counts), (2) flow-level statistics (bytes/second, packets/second, duration), (3) sub-flow and bulk statistics, and (4) TCP window and segment metrics. Features are extracted in the order defined by the CICIDS2017 column specification to ensure reproducible inference.

\textbf{Preprocessing}: NaN and infinite values are imputed with zero. Values exceeding the 99.9th percentile are clipped to prevent outlier-driven scaling distortion. For scale-sensitive models (Logistic Regression, SVM, MLP, KNN), a \texttt{StandardScaler} is fitted on the training set and applied to the test set. The fitted scaler is persisted alongside the model to ensure consistent transformation at inference time. A \texttt{LabelEncoder} maps the 15 class strings to integer indices 0--14.

\subsection{Model Selection}

Nine classifiers were evaluated: Logistic Regression (L2, balanced class weights), Decision Tree (GINI criterion, balanced weights), Random Forest (100 estimators, balanced weights), Extra Trees (100 estimators, balanced weights), AdaBoost (50 estimators), Gaussian Naive Bayes, K-Nearest Neighbors (k=5), Linear SVM (balanced weights), and MLP (two hidden layers: 100, 50 neurons). Gradient Boosting was excluded due to training time exceeding the 45-second sandbox constraint with 15 classes; XGBoost and LightGBM are recommended for production evaluation as noted in Section X.

Selection used a composite score balancing four metrics:

\begin{equation}
\text{Composite} = F1_{\text{macro}} \times 0.4 + \text{BalAcc} \times 0.3 + \text{MCC} \times 0.2 + \text{CV}_{F1} \times 0.1
\label{eq:composite}
\end{equation}

The weighting prioritizes macro F1 (equal per-class treatment) because raw accuracy is dominated by the BENIGN majority class. Balanced accuracy and MCC together penalize poor minority class performance. Cross-validation F1 measures generalization stability.

\subsection{Hyperparameter Tuning}

The top three baseline models (Decision Tree, Naive Bayes, Random Forest) underwent RandomizedSearchCV with 3-fold stratified cross-validation. Search spaces were:

\begin{itemize}
    \item \textbf{Decision Tree}: \texttt{max\_depth} $\in$ \{None, 10, 20, 30\}, \texttt{min\_samples\_split} $\in$ \{2, 5, 10\}, \texttt{criterion} $\in$ \{gini, entropy\}
    \item \textbf{Naive Bayes}: \texttt{var\_smoothing} $\in$ \texttt{logspace(-12, -6, 20)}
    \item \textbf{Random Forest}: Default \texttt{sklearn} configuration (\texttt{max\_features='sqrt'})
\end{itemize}

\subsection{Results}

Table~\ref{tab:baseline} reports baseline model performance. Table~\ref{tab:tuned} reports the tuned top-3. \textbf{Gaussian Naive Bayes (tuned)} achieves the highest composite score of 0.9731 with optimal parameters \texttt{var\_smoothing=1e-12}.

\begin{table}[h]
\caption{Baseline Model Comparison (Top 5 of 9)}
\label{tab:baseline}
\centering
\begin{tabular}{@{}lrrrrr@{}}
\toprule
\textbf{Model} & \textbf{F1} & \textbf{Acc} & \textbf{BalAcc} & \textbf{MCC} & \textbf{Comp.} \\
\midrule
Decision Tree        & 0.9967 & 0.9967 & 0.9542 & 0.9952 & 0.9774 \\
Naive Bayes          & 0.9952 & 0.9953 & 0.9189 & 0.9933 & 0.9563 \\
Random Forest        & 0.9931 & 0.9947 & 0.8489 & 0.9923 & 0.9361 \\
Extra Trees          & 0.9905 & 0.9923 & 0.8544 & 0.9890 & 0.8503 \\
Logistic Regression  & 0.9419 & 0.9383 & 0.7763 & 0.9122 & 0.7921 \\
\bottomrule
\end{tabular}
\end{table}

\begin{table}[h]
\caption{Tuned Model Comparison (Top 3 After RandomizedSearchCV)}
\label{tab:tuned}
\centering
\begin{tabular}{@{}lrrrrr@{}}
\toprule
\textbf{Model} & \textbf{F1} & \textbf{BalAcc} & \textbf{MCC} & \textbf{CV F1} & \textbf{Comp.} \\
\midrule
\textbf{Naive Bayes}  & \textbf{0.9692} & \textbf{0.9678} & 0.9957 & \textbf{0.9595} & \textbf{0.9731} \\
Decision Tree         & 0.9636 & 0.9600 & \textbf{0.9971} & 0.9454 & 0.9674 \\
Random Forest         & 0.8376 & 0.8489 & 0.9923 & 0.8577 & 0.8740 \\
\bottomrule
\end{tabular}
\end{table}

\subsection{Deployment Winner: Gaussian Naive Bayes}

Table~\ref{tab:winner} summarizes the deployed model. The selection rationale prioritizes the composite score, which Naive Bayes leads after tuning. Several operational properties reinforce this choice:

\begin{itemize}
    \item \textbf{Sub-millisecond inference}: GNB evaluates a 77-dimensional feature vector using 30 learned class mean/variance parameters in $O(n \cdot k)$ time, where $n=77$ features and $k=15$ classes. Measured inference latency is $<1$ ms/sample.
    \item \textbf{19 KB model footprint}: The serialized model, scaler, and label encoder occupy 19 KB combined, eliminating memory pressure in containerized deployments.
    \item \textbf{Strong minority class performance}: Balanced accuracy of 0.9678 indicates reliable classification of rare attack categories (Heartbleed, Infiltration) that would be suppressed by accuracy-optimizing models.
    \item \textbf{Stable generalization}: CV F1 of 0.9595 across 3 folds demonstrates low variance on unseen data splits.
\end{itemize}

The classifier is loaded once at startup and cached as a process-level singleton. A sparse input guard checks that at least 10 of 77 CICIDS2017 features are non-zero; when external events lack sufficient flow-level data, a heuristic rule-based fallback classifies based on destination port and flag counts with a fixed confidence of 0.65, ensuring all events receive at least a low-confidence classification.

\begin{table}[h]
\caption{Deployed Model Summary}
\label{tab:winner}
\centering
\begin{tabular}{@{}ll@{}}
\toprule
\textbf{Property} & \textbf{Value} \\
\midrule
Algorithm       & Gaussian Naive Bayes \\
\texttt{var\_smoothing} & $10^{-12}$ \\
Macro F1        & 0.9692 \\
Accuracy        & 0.9970 \\
Balanced Accuracy & 0.9678 \\
Matthews CC     & 0.9957 \\
CV F1 (3-fold)  & 0.9595 $\pm$ 0.000 \\
Composite Score & 0.9731 \\
Model version   & \texttt{v2.0.0-nb-tuned} \\
Model size      & $\approx$19 KB \\
Inference time  & $<$1 ms/sample \\
\bottomrule
\end{tabular}
\end{table}

% ─────────────────────────────────────────────────────────────────────────────
\section{Experimental Results}

\subsection{ML Classification Performance}

The deployed GNB classifier was evaluated on the held-out 3,000-sample test set. Table~\ref{tab:winner} reports the primary metrics. The near-perfect MCC of 0.9957 indicates strong discrimination across all 15 classes, including rare categories. The 3-fold CV F1 standard deviation of 0.000 (at 4 decimal precision) indicates exceptional stability, suggesting the model is not overfitting the training fold.

The composite scoring function in Equation~\ref{eq:composite} was validated by observing that the Decision Tree, while achieving higher raw accuracy (0.9967 vs 0.9970 for GNB, which appear similar but differ after post-tuning), achieves lower balanced accuracy (0.9600 vs 0.9678) indicating weaker minority class coverage — a critical failure mode in production intrusion detection where rare, high-severity attacks must not be missed.

\subsection{API Performance}

API latency measurements were conducted using the integration test suite against an SQLite in-memory database to isolate application logic from I/O. Key observations:

\begin{itemize}
    \item Authentication endpoint (\texttt{POST /api/v1/auth/login}): bcrypt verification is the dominant cost at approximately 100--200 ms on commodity hardware; this is an intentional security property of bcrypt \cite{provos1999bcrypt}.
    \item Incident list endpoint (\texttt{GET /api/v1/incidents}): sub-50 ms for paginated queries with indexed filters on PostgreSQL.
    \item ML classification inline with event ingestion: adds $<1$ ms per event, making it negligible relative to network and database I/O.
    \item PDF report generation: approximately 500--800 ms for a full incident report via reportlab, delivered as a streaming response.
\end{itemize}

\subsection{Evidence Integrity}

SHA-256 hash computation for evidence files is performed at upload time in Python using \texttt{hashlib.sha256}. The \texttt{/api/v1/evidence/\{id\}/verify} endpoint re-computes the hash of the stored file data and compares it to the original stored hash, returning a boolean \texttt{hash\_matched} field. Every access and verification event is recorded in the \texttt{chain\_of\_custody} table, creating an immutable audit trail. The \texttt{is\_immutable} flag on the \texttt{Evidence} model prevents programmatic deletion of forensic artifacts.

\subsection{Compliance Deadline Generation}

When an incident is created with jurisdiction flags (\texttt{personal\_data\_involved}, \texttt{health\_data\_involved}, \texttt{affected\_jurisdictions}), the Compliance Engine evaluates all three regulations and generates obligations for matching frameworks. Table~\ref{tab:compliance} summarizes the deadlines produced.

\begin{table}[h]
\caption{Compliance Obligation Generation Summary}
\label{tab:compliance}
\centering
\begin{tabular}{@{}lllr@{}}
\toprule
\textbf{Regulation} & \textbf{Jurisdiction} & \textbf{Authority} & \textbf{Deadline (hrs)} \\
\midrule
GDPR  & EU, EEA, UK & Data Protection Authority    & 72 \\
HIPAA & US          & HHS Office for Civil Rights  & 1,440 (60 days) \\
DPDPA & IN          & Data Protection Board India  & 72 \\
\bottomrule
\end{tabular}
\end{table}

Each regulation generates 3--4 specific obligation records with individual deadlines, enabling granular tracking of each notification step.

\subsection{Rate Limiting}

The sliding-window rate limiter maintains per-IP, per-path request counts in memory. Limits are: 10 requests/minute for \texttt{/api/v1/auth/login} and \texttt{/api/v1/auth/register}; 20 requests/minute for \texttt{/api/v1/auth/refresh}; 60 requests/minute for all other endpoints with a burst allowance of 20. Responses include \texttt{X-RateLimit-Limit} and \texttt{X-RateLimit-Remaining} headers. Exceeded limits return HTTP 429 with a \texttt{Retry-After} header.

% ─────────────────────────────────────────────────────────────────────────────
\section{Security Analysis}

\subsection{Authentication Security}

JWT tokens use HS256 signing with a minimum 32-character secret key enforced at startup. The \texttt{SECRET\_KEY} validator rejects the placeholder value in production environments. Access tokens carry a \texttt{jti} (JWT ID) unique claim; the logout endpoint records this claim in the \texttt{revoked\_tokens} table, enabling per-token blacklisting without session invalidation. Refresh tokens are long-lived (7 days) and issued only on successful login.

\subsection{Password Security}

Passwords are hashed using bcrypt with automatic work factor management via passlib. The bcrypt library version is pinned to 3.x (not 4.x) to avoid a known breaking change in password verification for passwords exceeding 72 bytes. Account lockout after 5 failed attempts for 15 minutes mitigates brute-force attacks at the application layer, complementing the rate limiter.

\subsection{Authorization Controls}

All protected endpoints resolve the caller's role from the JWT payload and evaluate the required permission against the centralized \texttt{ROLE\_PERMISSIONS} dictionary. Permission checks are the single source of truth; role string comparison is prohibited in application code. The \texttt{super\_admin} bypass for platform-level operations is always accompanied by an explicit audit log entry recording the bypassed resource.

\subsection{Security Headers}

The \texttt{SecurityHeadersMiddleware} appends security headers on every response: \texttt{X-Frame-Options: DENY} (clickjacking prevention), \texttt{X-Content-Type-Options: nosniff} (MIME-sniffing prevention), \texttt{X-XSS-Protection: 1; mode=block} (legacy XSS filter), \texttt{Referrer-Policy: strict-origin-when-cross-origin}, and \texttt{Permissions-Policy} disabling geolocation, microphone, and camera. HSTS (\texttt{max-age=31536000; includeSubDomains}) is applied on HTTPS connections. The server banner is removed from all responses.

\subsection{Project Isolation}

The multi-tenant data model enforces project boundaries at the ORM query layer. API key authentication resolves the project server-side; the client cannot submit or manipulate the project identifier. Every data query that returns incidents, evidence, or compliance records applies a \texttt{project\_id} filter. Cross-project access is structurally prevented rather than relying on application-layer checks.

\subsection{Evidence Immutability}

The \texttt{is\_immutable} field on the \texttt{Evidence} model defaults to \texttt{True}. The \texttt{EvidenceService.delete()} method explicitly rejects deletion of immutable evidence with a \texttt{PermissionDeniedError}. This preserves forensic artifacts from administrative deletion, supporting chain-of-custody requirements for legal proceedings \cite{casey2011digital}.

% ─────────────────────────────────────────────────────────────────────────────
\section{Compliance Framework}

\subsection{Design Philosophy}

The LBRO Compliance Engine follows a rule-trigger model: when an incident is created or updated with data classification flags (\texttt{personal\_data\_involved}, \texttt{health\_data\_involved}) and jurisdiction annotations (\texttt{affected\_jurisdictions}), the engine evaluates matching regulations and generates obligation records. This approach ensures compliance requirements surface immediately, eliminating the delay between breach discovery and legal awareness.

\subsection{GDPR}

GDPR Article 33 requires notification of the supervisory authority within 72 hours of becoming aware of a personal data breach. GDPR Article 34 requires notification of affected data subjects without undue delay if the breach is likely to result in high risk. LBRO generates four obligations per GDPR-matched incident: authority notification, subject notification, Article 33(5) register documentation, and risk assessment to natural persons. GDPR matching criteria: \texttt{personal\_data\_involved=True} or \texttt{affected\_jurisdictions} $\cap$ \{EU, EEA, UK\} $\neq \emptyset$.

\subsection{HIPAA}

HIPAA Breach Notification Rule (45 CFR \textsection164.400-414) requires covered entities to notify HHS within 60 days of discovery and affected individuals without unreasonable delay. Entities affecting more than 500 residents of a state must notify prominent media outlets. LBRO generates four HIPAA obligations and applies a 1,440-hour (60-day) deadline. HIPAA matching criteria: \texttt{health\_data\_involved=True} or jurisdiction \texttt{US}.

\subsection{DPDPA 2023}

India's Digital Personal Data Protection Act 2023 requires notification to the Data Protection Board within 72 hours of becoming aware of a personal data breach. LBRO generates three DPDPA obligations and applies a 72-hour deadline. DPDPA matching criteria: \texttt{personal\_data\_involved=True} or jurisdiction \texttt{IN}.

\subsection{Obligation Lifecycle}

Each \texttt{ComplianceRecord} transitions through three states: \textit{pending} (created, deadline not reached), \textit{overdue} (deadline past, not met), and \textit{met} (explicitly marked via \texttt{POST /api/v1/compliance/records/\{id\}/met} with supporting notes). The compliance dashboard aggregates records by regulation, showing total/met/overdue/pending counts and surfaces overdue items ordered by deadline ascending. Point-in-time compliance assessments are persisted as \texttt{ComplianceAssessment} records to support audit trails and compliance score trending.

\subsection{Limitations}

LBRO's compliance implementation covers the primary notification obligations of GDPR, HIPAA, and DPDPA. It does not yet implement: GDPR Data Protection Impact Assessment (DPIA) workflows; HIPAA Business Associate Agreement management; DPDPA consent management; PCI-DSS or SOC 2 frameworks; or jurisdiction-specific variations within the EU (e.g., Germany's stricter \texttt{BSI}-Gesetz requirements). The framework is designed for extension; adding a regulation requires a new entry in the \texttt{REGULATION\_RULES} dictionary with the applicable jurisdictions, deadline hours, authority, and obligation list.

% ─────────────────────────────────────────────────────────────────────────────
\section{Discussion}

\subsection{Strengths}

LBRO integrates components that are typically operated as separate products — detection (ML classifier), case management (incident workspace), digital forensics (evidence vault), and compliance (obligation engine) — into a single deployable stack. This integration eliminates inter-tool latency and ensures that every incident automatically receives forensic artifacts, compliance obligations, and investigation context at creation time.

The project-scoped API key ingestion model reduces integration friction to a single HTTP POST, making LBRO accessible to application developers without security operations expertise. The Demo Data Generator enables immediate evaluation of the full investigation workflow without requiring live security events.

The composite scoring function for model selection (Equation~\ref{eq:composite}) explicitly avoids accuracy-maximization traps common in class-imbalanced IDS benchmarks, prioritizing macro F1 and balanced accuracy to ensure rare, high-severity attack classes receive adequate classifier attention.

\subsection{Limitations}

Several honest limitations must be acknowledged:

\textbf{ML training data}: The deployed model was trained on a 15,000-sample synthetic subset of CICIDS2017. The full CICIDS2017 CSV ($\approx$2.8M flows) was not used due to sandbox constraints. GaussianNB assumes feature independence, which is violated by correlated CICIDS2017 flow statistics. Production deployment should evaluate XGBoost or LightGBM on the full dataset.

\textbf{Rate limiter persistence}: The current in-memory sliding-window rate limiter does not survive application restarts and does not share state across multiple API instances. Production horizontal scaling requires a Redis-backed implementation.

\textbf{Evidence storage}: File data is stored as PostgreSQL LargeBinary (bytea), which simplifies deployment but may create performance bottlenecks for large files or high-volume ingestion. The S3 service integration exists in the codebase but uses database storage as the primary path for reliability.

\textbf{Compliance interpretation}: The compliance engine generates obligations but does not provide legal advice. The obligation texts are informational; organizations must verify applicability with qualified legal counsel.

\subsection{Scalability}

The async FastAPI/asyncpg stack supports high concurrency on a single instance. The PostgreSQL connection pool (10 connections, 20 overflow) and SQS-backed worker queue provide horizontal scaling paths. The Terraform ECS deployment supports auto-scaling groups. The ML model's sub-millisecond inference adds negligible per-request overhead at any scale.

\subsection{Future Work}

Planned extensions include: XGBoost/LightGBM model evaluation on the full CICIDS2017 dataset; Redis-backed distributed rate limiting; federated compliance modules for PCI-DSS and SOC 2; a log agent (lbro-agent) for zero-code application monitoring; STIX/TAXII threat intelligence feed integration; and model drift monitoring with automated retraining triggers.

% ─────────────────────────────────────────────────────────────────────────────
\section{Conclusion}

This paper presented LBRO, a law-aware breach response orchestrator that integrates automated ML-based attack classification, tamper-evident digital forensics, and multi-jurisdiction regulatory compliance in a single open-architecture platform. A systematic evaluation of nine classifiers on CICIDS2017 identified Gaussian Naive Bayes (tuned, \texttt{var\_smoothing=$10^{-12}$}) as the optimal deployment choice with composite score 0.9731, macro F1 0.9692, balanced accuracy 0.9678, MCC 0.9957, and sub-millisecond inference latency with a 19 KB model footprint. The backend implements a four-tier RBAC system, per-token JWT revocation, bcrypt password security, sliding-window rate limiting, comprehensive security headers, and SHA-256 evidence integrity verification with immutable chain-of-custody tracking. The Compliance Engine automatically generates binding GDPR, HIPAA, and DPDPA obligations with deadline tracking at incident creation time, eliminating the gap between breach discovery and legal awareness. The complete stack is containerized with Docker Compose and deployable to AWS via Terraform. LBRO provides a reproducible, extensible foundation for incident response research and operational deployment, with particular applicability to SMEs and academic environments lacking dedicated security operations infrastructure.

% ─────────────────────────────────────────────────────────────────────────────
\section*{Acknowledgment}
The authors thank the Canadian Institute for Cybersecurity for the CICIDS2017 dataset and the open-source communities behind FastAPI, SQLAlchemy, scikit-learn, React, and the broader PyPI and npm ecosystems.

% ─────────────────────────────────────────────────────────────────────────────
\bibliographystyle{IEEEtran}
\bibliography{LBRO_IEEE_Paper}

\end{document}
